
\section{工作与科研经历}
\datedsubsection{\textbf{高级机器学习科学家, TikTok}}{2023/12 - 至今}
\begin{itemize}[parsep=0.5ex]
  \item 任职于TikTok（美国最大短视频平台）的电商业务算法团队，负责美国地区物流补贴算法从0到1的开发。
  \item 基于因果推断，设计并全量上线 Target Delivery Upgrade(TDU) 算法，通过个性化定制用户和商家的包邮门槛，将公司物流补贴ROI（投资回报率）提升16.7\%，销售额不变的情况，每天节约30万美元(经过A/B实验验证)。
  \item 连续两个半年绩效评定M+（在整个电商部门M-的背景下非常难得），获得2024H2 Spot Bonus (超出预期结果奖，4/100+)
\end{itemize}
\datedsubsection{\textbf{高级研究科学家, Lyft} }{2022/9 - 2023/12}
\begin{itemize}[parsep=0.5ex]
  \item 负责网约车司机奖励机制设计（Lyft核心团队之一），使用行为学、机器学习、优化等方法设计实时的“奖励地图”，目标是用最少的成本提高高峰期高需求地点的司机数量，本人担任“实时奖励地图”项目技术主管（Tech Lead）。
  \item 提出的新算法每年可为公司带来3800万美元的利润（通过A/B实验验证）。连续两次获得“Exceeding Expectation”绩效评定，仅用一年完成职级晋升，是Lyft公司历史上最快晋升（L4到L5）的校招博士生之一。
\end{itemize}
\datedsubsection{\textbf{算法科学家实习生, Lyft}}{2022/5 - 2022/8}
\begin{itemize}[parsep=0.5ex]
  \item 设计路径规划算法，用于共享自行车和电动滑板车的搬运，投放，更换电池等任务
  \item 提出了新型的车辆路径规划问题（VRP）的初始化方法，可提高 4.2\% 的运营效率，
每年可为公司节约 160 万美元的人力成本，目前已应用于公司世界范围内的所有市场
\end{itemize}
\datedsubsection{\textbf{研究助理, MIT} }{2018/8 - 2022/8}
\begin{itemize}[parsep=0.5ex]
  \item 在MIT Urban Mobility Lab和MIT Transit Lab研究与公共交通、需求与行为建模等相关的问题，参与多个基金项目申请，协助实验室组织工作（如团建、服务器管理，两次获得实验室乐于助人奖），协助导师指导MIT低年级研究生和本科生
\end{itemize}
\datedsubsection{\textbf{科研实习生, 芝加哥公共交通管理局（CTA）} }{2021/1 - 2021/2}
\begin{itemize}[parsep=0.5ex]
  \item 研究公共交通事故管理系统，开发了一套事故影响评估系统，可以迅速
分析和可视化公共交通事故影响。提出乘客路径推荐模型，可为芝加哥公共交通系统在事故期间降
低 20\% 的延误时间
\end{itemize}
\datedsubsection{\textbf{科研实习生, 香港地铁公司（MTR）} }{2019/7 - 2019/8}
\begin{itemize}[parsep=0.5ex]
  \item 开发了一套基于刷卡数据的自适应网络服务状态监控系统，被用于港铁日均 500 万人次出行的服务监
控，帮助港铁实时了解网络内的拥堵与延误信息，乘客的等待时间等。
\item 为港铁公司员工进行了四次培训讲座，协助他们学会使用该系统
\end{itemize}
\datedsubsection{\textbf{科研实习生, 新加坡-MIT联合科研中心（SMART）} }{2018/6 - 2018/8}
\begin{itemize}[parsep=0.5ex]
  \item 应用计量经济学与行为学模型研究人们对无人驾驶汽车的偏好。使用博弈论研究了无人驾驶汽车出现与公共交通的在市场上的竞争关系。该研究被 {\href{https://news.mit.edu/2021/smart-evaluates-competition-between-autonomous-vehicles-public-transit-0604}{\textcolor{blue}{{ MIT News}}}} 和多家其他媒体报道 (如： {\href{https://www.sciencedaily.com/releases/2021/05/210505102023.htm}{\textcolor{blue}{{ScienceDaily}}}}, {\href{https://autobala.com/smart-assesses-the-impact-of-competition-between-self-driving-cars-and-public-transport-sciencedaily/30616/}{\textcolor{blue}{{AutoBala}}}})
\end{itemize}
\datedsubsection{\textbf{研究助理, 清华大学土木工程系交通研究所} }{2015/12 - 2018/6}
\begin{itemize}[parsep=0.5ex]
  \item 与李瑞敏老师合作，使用车牌识别数据，进行车辆的速度曲线和尾气排放估计，以及动态OD出行矩阵估计
  \item 与李瑞敏老师和Rensselaer Polytechnic Institute的Cara Wang老师合作，研究了南宁市停车价格调整对停车需求和用户满意度的影响。
\end{itemize}
\section{教学经历}
\datedsubsection{\textbf{MIT本科生科研计划（UROP）}, }{2020/9 - 2022/5}
\begin{itemize}[parsep=0.5ex]
  \item \textbf{学术导师}
  \item 指导MIT本科生进行学术研究，主要研究公共交通事故对系统的需求和服务质量的影响。研究成果最终发表在IEEE Open Journal of Intelligent Transportation Systems
\end{itemize}
\datedsubsection{\textbf{波士顿全球教育公司（BGEC）}, }{2020/5 - 2022/5}
\begin{itemize}[parsep=0.5ex]
  \item \textbf{学术与科研项目主管}
  \item 设计公司学术与科研项目商业规划，设计并讲授了《机器学习与数据科学》、《机器学习：从基础到提高》两门课程，共同讲授《人工智能与机器人》课程。
\end{itemize}